\paragraph{On-Policy Distillation.} OPD computes supervision on trajectories sampled from the current student policy $\pi_\theta$. Given a prompt $x \sim \mathcal{D}_x$, the student samples a response sequence $\hat{y} = (\hat{y}_1, \dots, \hat{y}_N) \sim \pi_\theta(\cdot \mid x)$. Both the student model $\pi_\theta$ and the teacher model $\pi_T$ are evaluated on the student-generated prefixes $\hat{y}_{<t}$, yielding two next-token probability distributions at each step $t$: $p_t(v) \triangleq \pi_\theta(v \mid x, \hat{y}_{<t})$ and $q_t(v) \triangleq \pi_T(v \mid x, \hat{y}_{<t})$ over the vocabulary $v \in \mathcal{V}$.

A standard formulation of OPD minimizes the sequence-level reverse Kullback-Leibler (KL) divergence over the trajectories generated by the student:
\begin{equation}
\mathcal{L}_{\text{OPD}}(\theta) = \mathbb{E}_{x \sim \mathcal{D}_x} \Big[ D_{\text{KL}}(\pi_\theta(\cdot \mid x) \parallel \pi_T(\cdot \mid x)) \Big].
\end{equation}
Using autoregressive factorization, this sequence-level objective can be decomposed into an exact token-level formulation:
\begin{equation}
\mathcal{L}_{\text{OPD}}(\theta) = \mathbb{E}_{x \sim \mathcal{D}_x, \, \hat{y} \sim \pi_\theta(\cdot \mid x)} \left[ \sum_{t=1}^N D_{\text{KL}}(p_t \parallel q_t) \right].
\end{equation}

In practice, existing OPD formulations vary in their supervision granularity for computing this token-level KL, generally categorizing into sampled-token, full-vocabulary, and top-$k$ OPD. Among these, top-$k$ OPD restricts the divergence to a subset of high-probability tokens, significantly reducing memory overhead while preserving dense, multi-token supervision in the student's active generation region. Our experiments are primarily based on top-$k$ OPD.

\paragraph{Top-$k$ OPD.} Top-$k$ OPD restricts the divergence computation to a subset of tokens $S_t \subseteq \mathcal{V}$ at each step. Specifically, we adopt the student top-$k$ variant, which selects the $k$ tokens assigned the highest probabilities under the student distribution, defined as $S_t = \operatorname{TopK}(p_t, k)$. To compute the loss on this subset, we renormalize the student and teacher distributions over $S_t$:
\begin{equation}
\bar{p}_t^{(S_t)}(v) = \frac{p_t(v) \mathbf{1}[v \in S_t]}{\sum_{u \in S_t} p_t(u)}, \quad \bar{q}_t^{(S_t)}(v) = \frac{q_t(v) \mathbf{1}[v \in S_t]}{\sum_{u \in S_t} q_t(u)}.
\end{equation}
Distillation is then performed by minimizing the subset KL divergence $D_{\text{KL}}(\bar{p}_t^{(S_t)} \parallel \bar{q}_t^{(S_t)})$ over the trajectory, yielding the final objective:
\begin{equation}
\mathcal{L}_{\text{OPD}}^{\text{top-}k}(\theta) = \mathbb{E}_{x \sim \mathcal{D}_x, \, \hat{y} \sim \pi_\theta(\cdot \mid x)} \left[ \sum_{t=1}^N D_{\text{KL}}(\bar{p}_t^{(S_t)} \parallel \bar{q}_t^{(S_t)}) \right].
\end{equation}